\item \textbf{ACTIVIDAD: Localización del epicentro de un terremoto.}
El epicentro de un terremoto puede localizarse al observar la diferencia en los tiempos de llegada entre las ondas P y S en un sismógrafo. Una estación sismológica única puede determinar qué tan lejos se produce el terremoto. Con los datos de tres estaciones de sismógrafo, la triangulación se puede utilizar para determinar la ubicación exacta. Para los terremotos cercanos, las ondas sísmicas viajan a través de la corteza de la Tierra a velocidades típicas de 4.00 km/s para las ondas S y 8.00 km/s para las ondas P.

La siguiente tabla muestra las horas del día para la primera llegada de las ondas P y S en tres diferentes estaciones sismográficas de California:

\begin{center}
\begin{tabular}{|l|c|c|}
\hline
\textbf{Estación sísmica} & \textbf{Hora de llegada de la onda P} & \textbf{Hora de llegada de la onda S} \\ \hline
Sacramento & 3:21:34.4 PM & 3:22:04.8 PM \\ \hline
San Francisco & 3:21:35.9 PM & 3:22:07.8 PM \\ \hline
Los Ángeles & 3:21:52.1 PM & 3:22:40.2 PM \\ \hline
\end{tabular}
\end{center}

\begin{enumerate}
    \item Utilice los datos de la tabla para determinar qué gran ciudad de California representa el epicentro del terremoto.
    \item ¿A qué hora ocurrió el terremoto?
\end{enumerate}